% =========================================================================
%  Toward a Just Proposal:
%  Epistemic Injustice Theory and the Practice of Intimate-Relationship Proposals
%  Philosophy of Intimacy and the Theory of Justice — Paper VI
%  Wanhong Huang (Independent researcher)
%  Compile: xelatex main.tex && biber main && xelatex main.tex && xelatex main.tex
%
%  STATUS: theory-core draft. This installment sharpens the conceptual heart
%  before the full body is written: (i) a verified survey of the post-Fricker
%  literature and a detailed exposition of Catala (2025), and (ii) a mapping
%  with its necessary conditions and its boundaries against neighbouring concepts.
% =========================================================================
\documentclass[11pt,a4paper]{article}
\usepackage{serendip-paper}

% ── Bibliography (add refs.bib to this folder) ───────────────────────
\usepackage[backend=biber, style=authoryear, sorting=nyt, maxcitenames=2]{biblatex}
\addbibresource{../../bib/master.bib}

% Definition environment in series tone
\newenvironment{keydef}
  {\medskip\begin{quote}\color{ink}\itshape}
  {\end{quote}\medskip}

\begin{document}

% ── Title page ───────────────────────────────────────────────────────
\thispagestyle{empty}
\null\vfill
\begin{center}
{\footnotesize\color{rosedeep}\scshape Philosophy of Intimacy and the Theory of Justice \quad$\cdot$\quad Paper VI\par}
\vspace{0.9em}
\coverrule\par
\vspace{1.6em}
{\color{warnred}\Large Toward a Just Proposal: Epistemic Injustice Theory and the Practice of Intimate-Relationship Proposals\par}
\vspace{1.0em}
{\color{warnred}\large\itshape On the Marriage Proposal as a High-Stakes Decision, and the Conditions under Which Its Staging Wrongs the One It Asks\par}
\vspace{2.4em}
{\footnotesize\color{ink}\itshape Theory-core installment: literature verification and the framework it applies\par}
\vspace{2.0em}
\coverrule[0.25]\par
\vspace{1.4em}
{\large Wanhong Huang\par}
\vspace{0.5em}
{\footnotesize \texttt{huangwanhong@serendip.ngo}\par}
\vspace{1.2em}
{\today\par}
\end{center}
\vfill\null
\coverfootnote
\clearpage

% ── Dedication ───────────────────────────────────────────────────────
\clearpage
\thispagestyle{empty}
\vspace*{0.20\textheight}
\begin{center}
{\itshape\color{ink}
To the girl of the forest and the mountains,\\[0.4em]
to whom I would never put a question\\[0.25em]
she had not first been given every means to answer,\\[0.25em]
and whose yes, if it comes, I would have her free to have weighed.\\}
\vspace{2.6em}
{\cjkfont\large\color{rosedeep} 敬之惧之，匪以为难}\\[0.8em]
{\itshape\small\color{rosedeep}
To hold you in awe, to fear doing you wrong,\\[0.18em]
is not to make our love a hardship.}\\[1.0em]
{\footnotesize\color{lightgold}after the \cjkfont 诗经\/{\normalfont\ (\textit{Book of Songs})}}
\vspace{2.4em}

{\cjkfont\large\color{rosedeep} 信物者，待也，非约之封也}\\[0.8em]
{\itshape\small\color{rosedeep}
A token is a waiting,\\[0.18em]
not the sealing of a bargain.}\\[1.0em]
{\footnotesize\color{lightgold}a line of the author's own}
\end{center}
\vfill
\clearpage

% =========================================================================
\begin{abstract}
\noindent
A marriage proposal is, beneath its romance, a request that one person make a high-stakes, contract-like, life-structuring decision, and one that is, in L.~A.\ Paul's sense, \emph{transformative} \parencite{paul2014transformative}: its acceptance reshapes the values and self by whose lights it would be assessed. This paper argues that the conventional staging of the proposal, the engineered emotional peak, the social script that demands an immediate answer, and the asymmetry of information and preparation between the two parties, can, \emph{even in the complete absence of ill will and of any identity-based prejudice}, compromise the epistemic agency of the person asked.

\smallskip
\noindent \textbf{I do not claim to have discovered a new kind of epistemic injustice.} The wrong is better understood with a framework already in hand. I read the proposal through \textcite{catala2025dynamics}, \emph{The Dynamics of Epistemic Injustice: Situating Epistemic Power and Agency}, whose \emph{pluralist account of epistemic agency} and analysis of \emph{hermeneutical domination} and \emph{epistemic empowerment} supply exactly the apparatus the case requires. On Catala's account, epistemic agency is not exhausted by the propositional exchange of testimony; it includes a \emph{non-propositional, non-verbal, self-interpretive} dimension, the agent's capacity to make sense of her own experience to herself (self-to-self intelligibility) before and apart from conveying it to others (self-to-others intelligibility). The proposal, I argue, is a site where this self-interpretive agency, at the very threshold of a transformative choice, is placed under conditions, temporal, affective, informational, that impede its exercise.

\smallskip
\noindent This installment does not yet write the body. Following the principle that the conceptual core should be secured before the argument is built upon it, it does two things. \emph{First}, it verifies, against the literature current to 2025, the post-Fricker development of the field, Fricker's two originals, Dotson on testimonial smothering, Pohlhaus Jr.\ on willful hermeneutical ignorance, Medina on resistance, and the philosophical work on gaslighting, and then expounds Catala's framework in detail, since it is the framework the paper will use. \emph{Second}, it shows precisely how the proposal situation maps onto Catala's categories: which dimension of epistemic agency is at stake, why the wrong is a form of \emph{hermeneutical} rather than testimonial difficulty, and why locating it within an existing account is more defensible than minting a new term. The constructive half of the full paper will then derive the normativity of an \emph{epistemically just proposal}: separation of the affective moment from the decision, full disclosure before the asking, removal of immediate-answer pressure, the ring re-grounded from ``seal of a concluded bargain'' to ``token and waiting,'' and a deferred-commitment mechanism, each reconstrued as a condition of \emph{epistemic empowerment} in Catala's sense.

\vspace{0.6em}
\noindent\textbf{Keywords}\quad
epistemic injustice; epistemic agency; pluralist account; hermeneutical domination; epistemic empowerment; self-interpretation; transformative experience; marriage proposal; intimate justice; relational autonomy.
\end{abstract}

\tableofcontents
\clearpage

% =========================================================================
\section{Orientation: What This Installment Is and Is Not}\label{p06:sec:orientation}

This is the sixth paper in a series on the philosophy of intimacy and the theory of justice. The preceding papers established a relational ontology of the loving subject, the normativity of self-legislation and its re-grounding as co-legislation, and the standing of the other to contest the terms of a shared life. The present paper turns that apparatus on a single, concrete, and almost universally romanticised act: the marriage proposal.

The full paper will run a familiar two-part arc, a diagnosis followed by a constructive reconstruction. But the central claim, that the conventional proposal can wrong the person asked in her capacity as a knower and chooser, even when the proposer is loving, sincere, and wholly free of prejudice, lives or dies on the precision of one analysis. If that analysis is loosely drawn, the diagnosis collapses into either a familiar complaint (``surprises can be manipulative,'' which no one denies and which is not yet a theory) or an overreach (``all surprise proposals are unjust,'' which is false). The strategy of this series has been to let a neglected intimate situation \emph{illuminate} an existing theoretical structure rather than to multiply terminology.

I therefore make a deliberate methodological choice, and state it plainly because an earlier draft of this installment did the opposite. \textbf{I do not claim to have isolated a new species of epistemic injustice.} The temptation to coin one, ``situational'' or ``decisional'' epistemic injustice, was real, and an earlier version of this paper yielded to it. On reflection that move is both unnecessary and worse-supported than the alternative: the apparatus required to analyse the proposal already exists, in a precise and recent form, in \textcite{catala2025dynamics}. Minting a term where an adequate framework is in hand invites two avoidable costs, the suspicion of novelty-for-its-own-sake, and the burden of defending a neologism's boundaries from scratch. It is more defensible, and more illuminating, to show that Catala's pluralist account of epistemic agency, when trained on the proposal, \emph{already} names and explains the wrong. The contribution of this paper is then not a new concept but a new \emph{application}: bringing a frontier account of epistemic agency to bear on an intimate practice it was not written for, and deriving from that application a normative reconstruction of the practice.

So before any of the body is drafted, two preliminary tasks must be discharged with care.

\begin{enumerate}
\item \textbf{Verification of the field, and exposition of the framework.} Section~\ref{p06:sec:literature} reconstructs the post-Fricker terrain accurately, then expounds Catala's framework \emph{in detail}, because it is the framework the paper will use and the reader is owed more than its title.
\item \textbf{Mapping the case onto the framework.} Section~\ref{p06:sec:mapping} shows precisely how the proposal situation instantiates Catala's categories, which dimension of epistemic agency is engaged, why the difficulty is hermeneutical rather than testimonial, and where the analysis stops, so that the diagnosis is neither a truism nor an overreach.
\end{enumerate}

A methodological note on self-reference, discharged here so that it colours everything that follows. This paper, like its predecessors, is occasioned by the author's own practice; the constructive proposals it will defend are ones the author intends to enact. There is an obvious reflexive hazard: a paper that argues for the epistemic standing of one's intended partner, while using the shared relationship as its tacit case, would contradict its own thesis if it treated her as material to be theorised about rather than a person to be consulted. The series has handled this by anonymising and by presenting cases as constructed. The present paper accepts a stricter constraint, internal to its own argument: its constructive claims are not satisfied by being published but only by being \emph{enacted with} the other, disclosed to her, and held open to her revision, on pain of performative self-contradiction. I flag the hazard now and return to it as a structural feature, not an afterthought, in the full paper.

% =========================================================================
\section{The Field, Verified: Epistemic Injustice after Fricker}\label{p06:sec:literature}

The aim of this section is accuracy, not survey for its own sake. Each classical figure is reconstructed only far enough to locate the proposal case relative to it, and the reconstruction is checked against the literature current to 2025 rather than against memory; the section then turns to expound the framework the paper actually uses.

\subsection{Fricker's Two Forms, and Why Neither Fits}

\textcite{fricker2007epistemic} introduced \emph{epistemic injustice} as a wrong done to someone specifically in her capacity as a knower, and distinguished two species. \emph{Testimonial injustice} occurs when a hearer gives a speaker a deflated credibility owing to a prejudice attaching to the speaker's social identity, the paradigm being an \emph{identity-prejudicial credibility deficit}. \emph{Hermeneutical injustice} occurs when a gap in the collective interpretive resources leaves some group unable to render an important region of its experience intelligible, the standard example being the unavailability, before the term existed, of a concept like sexual harassment.

Two features of Fricker's originals matter here. First, both are tethered, in her formulation, to \emph{identity and structural marginalisation}: the credibility deficit tracks a prejudice against \emph{who the speaker is}, and the hermeneutical gap tracks the \emph{collective} exclusion of a marginalised group from meaning-making. Second, the wrong is located in a \emph{credibility} or \emph{intelligibility} deficit, in how a knower's contributions are received or her experience is conceptualised. The proposal case sits awkwardly with both. The person asked is not disbelieved because of her social identity; she is not, qua proposee, a member of a group excluded from interpretive resources; and the difficulty is not, in the first instance, that her contributions are received with a credibility deficit. What is at stake is closer to her capacity to make sense of her own situation \emph{to herself} under the conditions of the scene. Fricker's two originals therefore do not fit the case cleanly, not because the case is sui generis, but because Fricker's framework predates the pluralist widening of \emph{epistemic agency} that Catala supplies and that §\ref{p06:sec:catala} sets out. The right response is not a new label but the right framework.

\subsection{Dotson: Testimonial Smothering and Contributory Injustice}

\textcite{dotson2011tracking} identifies \emph{testimonial smothering}: a speaker truncates or withholds her own testimony because she anticipates that her audience will fail to receive it competently, so that the silencing is, in a sense, self-administered under structural pressure. \textcite{dotson2012cautionary} adds \emph{contributory injustice}, in which a knower's contributions are excluded because the audience wilfully relies on a biased set of interpretive resources while alternatives exist.

Dotson is the closest of the classical figures to the proposal case, and the comparison must therefore be made with precision rather than enthusiasm. Testimonial smothering shares with the proposal case the structural shape of a \emph{self-administered} epistemic shortfall: the proposee, like the smothered testifier, may pre-emptively suppress her own deliberation, here, the doubts and questions a momentous decision warrants, in anticipation of how raising them will be received. But the differences are decisive. In Dotson's smothering the content suppressed is \emph{testimony} (the speaker has knowledge she declines to deliver) and the anticipated failure is the audience's \emph{incompetence or hostility as a receiver}, rooted in structural marginalisation. In the proposal case what is suppressed is \emph{first-personal practical deliberation} (not knowledge to be delivered to another but reasoning the agent owes herself), and the pressure that suppresses it is not the proposer's incompetence as a receiver but the \emph{staging of the situation}, the public scene, the kneeling, the watching faces, the script that codes hesitation as cruelty. Dotson gives us the form ``epistemic shortfall produced under pressure rather than by an external withholder''; she does not give us the proposal's specific mechanism, which is situational design rather than anticipated mis-reception, and which need not trade on any structural marginalisation of the proposee at all.

\subsection{Pohlhaus Jr.: Willful Hermeneutical Ignorance}

\textcite{pohlhaus2012relational} develops \emph{willful hermeneutical ignorance}: dominantly situated knowers actively refuse to take up the interpretive resources developed by the marginalised, so that the marginalised are met with a wilful failure to understand. This is, again, a concept built on the axis of \emph{structural power between social positions} and on the fate of \emph{interpretive resources}. It is a useful marker of an edge, not a home for the case. The proposal difficulty, as analysed below, does not require that the proposer occupy a dominant social position relative to the proposee, nor that any interpretive resource be refused; it can arise between social equals, and indeed between partners where the proposee is in every structural respect the proposer's equal or superior. Pohlhaus marks one edge of the territory; the home, as the next subsection argues, is Catala's pluralist account of agency.

\subsection{Medina: The Epistemology of Resistance}

\textcite{medina2013epistemology} shifts the register from individual acts to \emph{epistemic character, responsibility, and resistance}: he theorises the epistemic vices of the privileged (arrogance, laziness, closed-mindedness), the virtues that resist them, and the way epistemic injustices are sustained or contested across a social ecology. Medina is less a boundary than a \emph{resource for the constructive half}. The notion of an epistemically just proposal that the full paper will defend is, in Medina's terms, a structure of \emph{epistemic responsibility} on the proposer's side, an active cultivation of the conditions under which the other can exercise full epistemic agency, and of \emph{epistemic resistance} as something the staging should make available rather than foreclose. I record him here so that the constructive section can draw on him without re-introducing him cold.

\subsection{Gaslighting: The Indispensable Contrast}

The philosophical literature on gaslighting is the contrast that most sharpens the proposal case, because gaslighting is the canonical wrong to epistemic agency \emph{within an intimate relationship}, and the force of the present analysis is that the proposal can compromise epistemic agency \emph{without} the feature that literature treats as central.

The state of that literature, verified, is as follows. \textcite{abramson2014turning} gives the influential account on which gaslighting is constitutively a matter of the gaslighter's \emph{intention and motive}, an attempt to gain control by destroying the target's standing as an independent locus of judgement, and on which its wrong is primarily moral rather than epistemic. \textcite{spear2019epistemic} and the surrounding work argue, against Abramson's de-emphasis of the epistemic, that gaslighting has ineliminable epistemic dimensions, that it corrodes \emph{self-trust} and the victim's conception of herself as an autonomous knower. \textcite{stark2019gaslighting} extends gaslighting beyond the dyadic and intentional toward the structural, tying it to prior epistemic injustice. \textcite{podosky2021gaslighting} argues that gaslighting can be understood as a \emph{purely epistemic} phenomenon, characterised by confronting the victim with a forced choice between rejecting the gaslighter's testimony and doubting her own basic epistemic competence, and need not depend on a particular set of intentions; \textcite{ruiz2020cultural} theorises \emph{cultural} gaslighting at the level of colonial and structural epistemic oppression.

Two things follow. First, even on the most deflated, least intention-dependent account (Podosky), gaslighting is defined by an attack on the victim's \emph{self-trust or basic epistemic competence}, by getting her to doubt that she knows what she knows. The proposal case leaves the proposee's self-trust and competence \emph{intact}: she is not made to doubt her faculties; she is placed in a situation in which those intact faculties cannot be fully exercised before she must answer. This is a difference in the \emph{locus} of the difficulty, competence undermined versus exercise impeded, and it is what tells us we are not looking at gaslighting and should not analyse the case as a covert or attenuated form of it. Second, every account in the gaslighting literature, even the structural ones, treats the wrong as flowing from something the wronging party \emph{does to} the victim's epistemic self-relation, whether by intention (Abramson), by epistemic mechanism (Podosky), or by structural position (Stark, Ruíz). The proposal difficulty can flow from a \emph{situation the proposer has innocently inherited and reproduced}, the cultural script of the romantic proposal, without his doing anything to the proposee's self-relation and without his intending or benefiting from any control. Gaslighting is thus the contrast that tells us \emph{where not to look}: not to malice, not to a damaged self-relation. It does not, by itself, tell us where to look instead. For that we need Catala.

\subsection{Catala's Framework, Expounded}\label{p06:sec:catala}

The account this paper uses is \textcite{catala2025dynamics}, \emph{The Dynamics of Epistemic Injustice: Situating Epistemic Power and Agency} (Oxford University Press, 2025). Because the paper rests on it, more than a title is owed; this subsection sets out the parts of the framework the argument will use. (The exposition is faithful to the book's published structure and stated commitments; particular formulations the full paper will source to specific chapters once the page references are settled.)

\subsubsection{The two undertheorised notions, and the method}

Catala's animating observation is that the literature on epistemic injustice constantly invokes \emph{epistemic power} and \emph{epistemic agency} yet leaves both largely undertheorised. Her project is to give a systematic account of the two by working through the \emph{dynamics} of epistemic injustice, the many forms it takes, the various sites and mechanisms through which it operates, and the transformations required to cultivate epistemic justice. Methodologically she adopts standpoint theory both as theory and as method: epistemic power and agency are to be understood as \emph{situated}, indexed to social location, rather than as the properties of an abstract knower. Two of her headline contributions matter here: a systematic account of epistemic power and agency that foregrounds the \emph{interaction between individual and structural factors} (neither pure individualism nor pure structuralism), and a \emph{pluralist} account that brings into view the \emph{non-propositional and non-verbal} dimensions of epistemic agency that earlier, testimony-centred accounts neglected.

\subsubsection{The pluralist account of epistemic agency}

This is the part the proposal case turns on. For Fricker the paradigm of epistemic agency is \emph{testimonial}: a speaker conveys a proposition to a hearer who assigns it some credibility. Catala's pluralism widens the field. Epistemic agency, on her account, is the capacity to be an actor in one's activity of knowing, to produce, use, or convey knowledge, and this capacity has dimensions that are not propositional and not verbal at all. Centrally for us, it includes a \emph{self-interpretive} dimension: the agent's capacity to render her own experience intelligible \emph{to herself}, prior to and independently of any conveying of it to another. Following the distinction the framework draws, intelligibility runs along two axes, \emph{self-to-others} (being understood by an interlocutor) and \emph{self-to-self} (understanding one's own situation). A deficit on the second axis is a failure or foreclosure of self-interpretation: the agent cannot, in the conditions she is in, make adequate sense of her own experience to herself. This is a genuinely \emph{hermeneutical} matter, but a hermeneutical matter relocated, at least in part, \emph{inside} the agent's relation to her own experience rather than residing wholly in the collective pool of concepts.

\subsubsection{Testimonial and hermeneutical domination}

Catala reframes the two Frickerian species in terms of \emph{domination}, the better to capture the interaction of individual and structural factors. \emph{Testimonial domination} (her Chapter~2) gives a structural explanation of how stereotypes systematically deflate the epistemic agency of the dominated as conveyers of knowledge. \emph{Hermeneutical domination} (her Chapter~3) concerns the conditions under which the dominated are impeded in making sense, whether to others or to themselves, with phenomena such as ``white ignoring'' and deliberative impasse as mechanisms. The vocabulary of \emph{domination} (rather than discrete ``injustice events'') is deliberate: it lets the account track the standing structural conditions that bear on the exercise of agency, not only datable acts. For the proposal, the relevant register is \emph{hermeneutical}: what is at issue is the conditions of sense-making, not the assignment of credibility to testimony.

\subsubsection{Becoming who you are: transformative experience and epistemic empowerment}

Catala's Chapter~6, \emph{Becoming Who You Are: Hermeneutical Breakthroughs, Transformative Experience, and Epistemic Empowerment}, is the hinge between the diagnostic and the constructive use this paper will make of her. It connects epistemic agency to \emph{transformative experience} in roughly L.~A.\ Paul's sense, experiences whose undergoing reshapes the values and the self by whose lights one would have deliberated about whether to undergo them, and theorises \emph{hermeneutical breakthrough} and \emph{epistemic empowerment} as the positive achievements by which an agent comes to make new sense of herself across such a threshold. Where hermeneutical domination names the impeding of self-interpretation, epistemic empowerment names its enabling and enlargement. This gives the paper its constructive vocabulary: an epistemically just proposal is one structured to \emph{empower} the proposee's self-interpretation at the threshold of a transformative choice, rather than to compress or crowd it.

\subsubsection{Why this framework, and an honest note on fit}

Two honest qualifications. First, Catala's sustained cases concern \emph{non-dominant social groups} (academic migrants, intellectually disabled persons, autistic women, divided societies); the proposee in a loving, equal partnership is typically not dominated in that structural sense. The paper therefore uses Catala's \emph{analytic apparatus}, the pluralist account of agency, the self-to-self axis of intelligibility, the hermeneutical register, the empowerment/transformation pairing, without claiming the proposee is a member of a structurally marginalised group. This is a legitimate use of a framework: borrowing its machinery while being explicit about which of its motivating conditions do and do not obtain. Second, precisely because the structural-domination reading does not straightforwardly apply, the paper must argue that the \emph{situational} conditions of the proposal (temporal, affective, informational) can impede self-interpretive agency even absent group-based domination, an extension of the framework that §\ref{p06:sec:mapping} begins and the full paper defends. Naming this honestly is better than smuggling the proposee into a category of marginalisation she does not occupy.

% =========================================================================
\section{The Jurisprudence of the Proposal: A Quasi-Contractual Act under Competing Theories of Justice}\label{p06:sec:jurisprudence}

The epistemic analysis above asks whether the proposee's self-interpretive agency is empowered or impeded. A jurisprudential analysis asks a different and complementary question: treating the proposal as what it structurally is, the solicitation of a momentous, quasi-contractual act of consent, \emph{is the act just}, and what would each major theory of justice say about the conditions under which it is or is not? This section runs that analysis. It is complementary rather than redundant: the epistemic register concerns the \emph{knower}; the jurisprudential register concerns the \emph{legitimacy of the consent} and the \emph{order it institutes}. The series has used this method before (the Hohfeldian and republican analysis of the vow in Paper~III); here it is trained on the proposal as a consent-act.

A caveat fixes the scope. To call the proposal ``quasi-contractual'' is an analytic device, not a reduction of love to bargaining. A marriage proposal is not an offer in contract law, and its acceptance is not consideration. But it shares with consent-acts the features that the theories of justice were built to assess: it solicits, at a moment, a binding answer that institutes an order of mutual claims and obligations and that is costly to revise. That shared structure is what licenses the analysis; the disanalogies are flagged where they bear.

\subsection{The Consent-Theoretic Frame: Four Conditions of Valid Consent}\label{p06:sec:consent}

Before the schools diverge, they share a common analytic inheritance from the theory of consent, refined most precisely in the bioethics and research-ethics literatures. On the standard account, valid consent to a serious intervention requires four things: \emph{capacity} (the ability to engage in reasoned deliberation), \emph{disclosure} (the sharing of information material to the decision), \emph{understanding} (an actual appreciation of the nature, risks, and alternatives), and \emph{voluntariness} (freedom from coercion and undue influence) \parencite{eyal2019informed, beauchamp2019principles}. Consent is compromised when any element is lacking.

Two refinements from that literature bear directly on the proposal. First, the standard for voluntariness and understanding is not fixed but \emph{scales with the stakes}: the graver and less reversible the decision, the greater the clarity, authenticity, coherence, and settled commitment that valid consent requires \parencite{roberts2002voluntarism}. A binding answer to a life-structuring, transformative question therefore demands \emph{more}, not less, of the four conditions than a casual agreement does. Second, voluntariness is undermined not only by literal coercion (a threat that would leave one seriously worse off) but by \emph{undue influence} and by ``no choice'' framings that constrict the practical space of refusal \parencite{eyal2019informed}.

Mapped onto the conventional staged proposal, the diagnosis is immediate and theory-neutral, which is its strength. \emph{Capacity} is intact (this is not gaslighting). But \emph{disclosure} is impaired by informational asymmetry; \emph{understanding} is impaired by affective saturation crowding the interpretive room; and \emph{voluntariness} is impaired by the immediate-answer script, which functions as a soft ``no choice'' framing, where public refusal carries a staged social cost. By the very standard that scales \emph{up} with the stakes, the conventional proposal solicits consent under conditions that the consent literature would flag as falling short, even with a loving, sincere, non-coercive proposer. The constructive remedies of the full paper (prior disclosure, decoupled timing, removed immediacy) are, in this register, simply the conditions of \emph{valid consent to a high-stakes act}. This is the shared floor; the schools below differ on what justice additionally requires.

\subsection{Liberalism I: Rawlsian and Procedural Justice}

A Rawlsian analysis does not directly govern the intimate dyad (justice as fairness regulates the basic structure, not the proposal), but its \emph{method} transfers. Ask what procedure for soliciting a momentous commitment the two parties would agree to behind a veil of ignorance, not knowing which role, proposer or proposee, they would occupy. No one ignorant of her position would choose a procedure that loads information, timing, and emotional framing onto one side and demands the other answer cold and at once; each would choose the procedure that protects the asked party, since either might be that party. Procedural justice thus condemns the staged proposal not for its outcome but for its \emph{procedure}: fair terms of agreement are those the disadvantaged party would accept in advance, and the immediate-answer-under-asymmetry procedure fails that test. The Rawlsian verdict: \emph{the conventional proposal is procedurally unjust; the just proposal is the one whose procedure both would choose without knowing who would be asked}.

\subsection{Liberalism II: Autonomy-Based and Consent Theories}

The autonomy-centred liberal (Kantian or Millian) locates justice in respect for the agent's self-governing choice. Here the consent frame of §\ref{p06:sec:consent} does the work directly: to solicit a binding answer under impaired voluntariness and understanding is to fail to treat the proposee as a fully self-governing end, to take an answer that does not fully express her autonomous will. The Kantian formulation is sharp: staging that uses affect and time-pressure to elicit a ``yes'' treats the proposee's emotional response as a \emph{means} to securing agreement, in tension with the requirement to treat her rational agency as an end. The verdict: \emph{the conventional proposal risks using the proposee's own affect against her autonomy; the just proposal is one to which her yes is the act of a self-governing will, fully informed and unhurried}.

\subsection{Republicanism: Non-Domination and Relational Autonomy}

This is the register of Paper~III, carried forward. The republican does not ask whether interference occurred but whether one party holds \emph{arbitrary power} over the other. The staged proposal, considered alone, is not yet domination, a single act, not a standing power. But it can be the \emph{instituting moment} of an order in which second-order power (the power to set and revise the terms of the shared life) is unequally distributed, and a proposal whose form already places one party in the position of reactive, time-pressured respondent can inaugurate that asymmetry. The republican adds what consent theory omits: that what matters is not only the validity of this one answer but the \emph{standing} it institutes, whether the proposee retains live, robust power to contest the terms thereafter. The verdict: \emph{the proposal is just only if it institutes a relationship in which the asked party retains undominated standing to revisit and revise; a proposal that secures a ``yes'' precisely by foreclosing deliberation is a poor omen for, and may be the first act of, an order of benevolent domination}.

\subsection{Marxian and Social-Reproduction Theory}

The Marxian register, also carried from Paper~III, reads the proposal as the entry-point to an institution, marriage, historically organised around the appropriation of unwaged relational and reproductive labour. The proposal's \emph{form} is then significant: a staging that secures rapid agreement under emotional and social pressure, before the material terms of the shared life (who will carry the domestic and care labour, whose career bends to whose) are disclosed and negotiated, functions ideologically. It recodes the entry into a labour-asymmetric institution as a moment of pure romance, so that the unequal terms are agreed to before they are seen. The verdict: \emph{the conventional proposal can be an ideological act that obtains agreement to an exploitative division of labour by ensuring it is never on the table at the moment of consent; justice requires that the material terms of social reproduction be disclosed and open to negotiation before, not after, the binding answer}.

\subsection{Feminist Justice I: Pateman and the Critique of the Contract Itself}\label{p06:sec:pateman}

The deepest challenge comes from feminist contract critique, and the paper must meet it rather than evade it, because it threatens the constructive project at its root. \textcite{pateman1988sexual} argues that the contractual frame cannot be reformed into justice: contract, in her analysis, always generates political right in the form of domination and subordination, and the marriage contract specifically rests on a presumed ``law of male sex-right'' that the language of free agreement masks. Crucially for this paper, Pateman explicitly rejects the move this paper is tempted by, the hope that a contract ``between genuinely equal partners'' or ``entered into without any coercion'' could be redeemed. On her view, the feminist who seeks a \emph{properly} consensual marriage contract is ``misleading'' herself: the defect is structural, not procedural, and lies in the very figure of the person as proprietor contracting over herself.

The honest response is not to refute Pateman but to locate the paper's claim precisely against her. Three points. First, the paper does not claim that perfecting the consent-conditions of the proposal \emph{achieves} justice in marriage; it claims, more modestly, that impaired consent at the threshold is \emph{one} remediable wrong, whose remedy is necessary though not sufficient. Pateman's structural critique and this paper's situational one can both be true. Second, Pateman's target is ``property in the person'' as the hidden premise of the contract; the relational reconstruction of Papers~III and~V already rejects that premise, the vow re-grounded as co-legislation is precisely \emph{not} a transfer of property in the person but an ongoing shared authorship. To the extent the proposal is reconceived as the opening of co-legislation rather than the sealing of a transfer, it steps off the contractual terrain Pateman condemns, which is also why the ring is re-grounded from ``seal of a bargain'' to ``token and waiting.'' Third, where Pateman is right that no staging can dissolve the historical-structural asymmetry of the institution, the paper concedes the residue openly: the just proposal is the most one can do \emph{at the threshold}, and it neither claims nor pretends to undo what only the transformation of the institution could undo. The verdict: \emph{Pateman shows that procedural justice at the proposal is not sufficient for justice in marriage; she does not show it is not necessary, and the relational re-grounding answers her premise of property-in-the-person without claiming to have answered her structural critique whole}.

\subsection{Feminist Justice II: Relational Autonomy and the Ethics of Care}

A second feminist tradition, less sceptical of reform, supplies positive standards. The relational-autonomy theorists hold that autonomy is not the property of an isolated chooser but is constituted and sustained in relationships; a choice is autonomous to the degree the relationship \emph{supports} the agent's self-governing capacities rather than undercutting them \parencite{mackenzie2000relational}. On this view a proposal is just when its form actively scaffolds the proposee's reflective agency, by giving time, information, and room for her own interpretation, rather than relying on her competence surviving an unfavourable staging. The ethics of care adds that the relevant moral salience is \emph{responsiveness to the concrete other}: a proposer attentive to this particular person's needs at the threshold would not stage a scene optimised for his narrative of the romantic moment, but one optimised for her capacity to answer well. The verdict, convergent with this paper's constructive half: \emph{justice in the proposal is the active, caring construction of the conditions under which the other can exercise full relational autonomy, the proposal as a gift of deliberative room, not a demand for immediate ratification}.

\subsection{Virtue, Care, and the Confucian \texorpdfstring{\cjkfont 礼义}{liyi} Register}

A final register, drawing the Aristotelian, care-ethical, and Confucian threads together, shifts from the justice of the act to the \emph{character} it expresses and the \emph{ritual propriety} (\cjkfont 礼\,\emph{l\v{i}}) through which care is rightly enacted. The Confucian tradition is distinctively apt here: it does not oppose emotion and propriety but holds that right feeling is \emph{realised} through proper form, ``\cjkfont 发乎情，止乎礼义'', it rises from feeling and comes to rest in what is right (the line that stands in this series' dedication). On this view the staged surprise proposal errs not because it is emotional but because its form is \emph{self-regarding}: it expresses the proposer's feeling in a manner that subordinates the beloved's standing to answer, where \cjkfont 礼 would shape the same feeling into a form that honours her, that makes space, gives notice, awaits readiness. The virtue at stake is a species of \emph{respect-in-love} (\cjkfont 敬), the holding of the beloved in a regard that fears to do her wrong (\cjkfont 敬之惧之). The verdict: \emph{the just proposal is the one whose form is the outward propriety (\cjkfont 礼) adequate to the inward love (\cjkfont 情), in which awaiting the other's readiness is not a diminishment of ardour but its disciplined and respectful expression}.

\subsection{Convergence, Divergence, and What the Body Owes}

The schools converge more than they diverge, which is itself a result worth stating: liberal procedural justice, autonomy-based liberalism, republicanism, relational autonomy, the ethics of care, and the Confucian \cjkfont 礼义 register all condemn the conventional staged proposal and all endorse, by different routes, the same family of remedies, disclosure, decoupled timing, removed immediacy, retained standing to revise. The Marxian register adds a condition the others underweight (disclosure of the \emph{material} terms of social reproduction). And Pateman dissents at the deepest level, denying that any procedural fix redeems the contractual form, a dissent the paper answers only partially and concedes openly. The body of the paper will (i) develop the consent-theoretic diagnosis of §\ref{p06:sec:consent} in full, (ii) carry the republican and Marxian lines from Paper~III without simply repeating them, (iii) treat Pateman's objection as the principal challenge to the constructive project and give §\ref{p06:sec:pateman}'s three-point reply its full defence, and (iv) show that the convergence of the remaining schools on the same remedies is not coincidence but reflects that those remedies just \emph{are} the conditions of valid, undominated, caring consent to a transformative commitment.

We can now show how the proposal situation instantiates Catala's categories, without minting a new term. I proceed in four steps: which dimension of agency is engaged; why the register is hermeneutical and not testimonial; how the situational conditions impede self-interpretation; and where the analysis stops.

\subsection{The Proposal as a Transformative Threshold}\label{p06:sec:mapping}

Marriage is a paradigm of transformative experience in L.~A.\ Paul's sense: one cannot fully know, in advance, what it will be like to be married to this person, nor who one will become in the marrying, and so the values by which one would assess the choice are partly constituted by undergoing it. Catala's Chapter~6 is built for exactly this juncture: it theorises how an agent makes new sense of herself, achieves a \emph{hermeneutical breakthrough}, across such a threshold, and what it is for her to be \emph{epistemically empowered} rather than impeded in doing so. The proposal is the moment at which a person is asked to commit, by a single binding answer, across precisely such a threshold. The question the framework lets us ask is sharp: \emph{under what conditions is the proposee's self-interpretive agency, at this threshold, empowered, and under what conditions impeded?}

\subsection{Which Dimension of Agency: The Self-to-Self Axis}

The wrong, if there is one, is not that the proposee is disbelieved (no credibility is at issue) and not that she lacks the public concept of marriage (the collective hermeneutical resource is amply available). It is that, in the staged scene, she is impeded in the \emph{self-to-self} exercise of intelligibility: making adequate sense, to herself, of what this commitment means for her, what she wants, what she fears, who she would become, in the moment she must answer. This is squarely the non-propositional, self-interpretive dimension of epistemic agency that Catala's \emph{pluralist account} brings into view and that testimony-centred accounts miss. Naming the dimension precisely is what keeps the diagnosis from collapsing into the truism that ``surprises can be manipulative'': the claim is specific, that a particular and theorised capacity, self-to-self intelligibility at a transformative threshold, is the thing the staging bears on.

\subsection{Why Hermeneutical, Not Testimonial}

Catala's reframing as \emph{domination} along testimonial and hermeneutical axes lets us place the case exactly. It is not \emph{testimonial}: no one is deflating the credibility of the proposee's word. It belongs to the \emph{hermeneutical} register, the conditions of sense-making, but with a twist the framework is built to accommodate: the impediment is not (as in Fricker's original) a gap in the collective conceptual pool, but a situational compression of the agent's \emph{occasion and capacity to interpret her own experience to herself}. Catala's self-to-self axis is what makes this sayable without distortion. The proposal does not deprive the proposee of the concept of marriage; it deprives her, in the scene, of the conditions under which she could bring her own interpretive agency to bear on her own case.

\subsection{How the Situational Conditions Impede Self-Interpretation}

Three features of the conventional staging bear on self-to-self intelligibility, and it is worth being exact about the mechanism in each, since the constructive half will work by reversing them.

\begin{enumerate}
\item \textbf{Affective saturation.} A staged emotional peak does not merely accompany the decision; it can crowd the cognitive and interpretive room in which the proposee would make sense of her own response, so that what presents itself as clarity may be the affect of the scene rather than a self-interpretation she has had the room to reach. Emotion is not the enemy, it is proper to love, but emotion \emph{staged to displace interpretation at the moment interpretation is most needed} is an impediment to self-to-self intelligibility.
\item \textbf{Temporal compression and the immediate-answer script.} The script codes hesitation as rejection or cruelty and demands an answer \emph{now}, collapsing the interval in which self-interpretation could occur and converting the need to think into a social cost borne in public. This is a direct constriction of the occasion for hermeneutical agency.
\item \textbf{Informational asymmetry.} One party has had unbounded time to interpret and prepare; the other meets the threshold cold. The asymmetry is not (or not only) about propositional facts withheld; it is about the unequal opportunity for the \emph{interpretive work} that a transformative choice demands.
\end{enumerate}

Each of these, in Catala's terms, bears on \emph{epistemic empowerment}: the conventional staging tends to \emph{disempower} the proposee's self-interpretation at the threshold, where the normative ideal would be to empower it.

\subsection{Where the Analysis Stops: Two Honest Limits}\label{p06:sec:limits}

Two limits keep the application from over-reaching, and both follow from taking Catala seriously rather than from an ad hoc fence.

First, \emph{the structural-domination gap}. Catala's domination concepts are indexed to non-dominant social groups; the proposee in a loving, equal partnership is not, qua proposee, structurally marginalised. The paper therefore claims only that the \emph{situational} conditions can impede self-interpretive agency, and does not claim the proposee suffers \emph{hermeneutical domination} in Catala's full structural sense. This is the extension the framework needs and the full paper must defend; it is named, not hidden. Where a real power asymmetry \emph{does} obtain between the parties, the structural reading applies more directly, and the wrong is correspondingly graver.

Second, \emph{antecedent interpretation dissolves the case}. If the couple has, over months, done the interpretive and deliberative work together, then a staged, surprise, immediate-answer proposal compresses no self-interpretation that has not already had its occasion; the answer ratifies a long joint sense-making rather than substituting for it. This is exactly why the constructive remedy is \emph{prior disclosure and shared interpretation}: dissolve the asymmetry and restore the occasion in advance, and the same staging that would otherwise impede self-to-self intelligibility now impedes nothing. The analysis points directly at its own remedy, which is the mark of a diagnosis that has located the wrong in the right place.

\subsection{What the Body Will Build on This}

With the framework fixed and the mapping done, the full paper's arc is set. The \emph{diagnosis} applies §\ref{p06:sec:mapping} to the conventional proposal: in the sincere, prejudice-free, malice-free case, the staging can still impede the proposee's self-to-self intelligibility at a transformative threshold, an impediment to epistemic agency in Catala's pluralist sense, and no appeal to a new category is required to say so. Three objections are then met: the \emph{kill-joy} objection (refusing to let affect displace interpretation is not refusing affect); the \emph{voluntary-participation} objection (participation under a compressed occasion for self-interpretation is not the same as having had that occasion); and the \emph{scope} objection (the target is the structure of soliciting a transformative commitment under asymmetry and immediate-answer pressure, not every surprise). The \emph{constructive} half then derives the epistemically just proposal as a structure of \emph{epistemic empowerment}: separation of the affective moment from the decision; full disclosure and shared interpretation before the asking; removal of immediate-answer pressure; the ring re-grounded from \emph{seal of a concluded bargain} to \emph{token and waiting}; and a deferred-commitment mechanism. The reflexive constraint of §\ref{p06:sec:orientation} governs throughout: the argument is honoured only if enacted \emph{with} the other, not merely asserted about her, which is itself an application of empowering, rather than impeding, her self-interpretation.

% =========================================================================
\section*{Acknowledgements}
\addcontentsline{toc}{section}{Acknowledgements}

This is the sixth paper in a series on the philosophy of intimacy and the theory of justice. It was occasioned, like its predecessors, by a question the author's own life put to him; the proposee discussed here is kept deliberately anonymous, and the case is presented as a constructed one, so that the argument may stand on its structure. The reflexive constraint announced in the orientation is meant sincerely: the paper's claims are not discharged by publication but only by being enacted, disclosed, and held open to revision with the one they concern.

In the interest of transparency, I note that an AI assistant was used in preparing this manuscript, as a tool for verifying the current literature, structuring the analysis, and refining the prose; the ideas, commitments, and final judgements are my own, and I take full responsibility for the content. The literature of this installment was verified against sources current to 2025; any remaining error of attribution is mine to correct in the full draft.

\begin{center}
{\cjkfont\color{rosedeep} 愿天下有情人，问者以诚，答者以宁；相知而后相许，从容而不相负。相许如初见，白首如相知。}
\end{center}

% =========================================================================
\printbibliography[heading=bibintoc]

\end{document}
